\DocumentMetadata{
  pdfversion=2.0,
  pdfstandard=ua-2,
  lang=en-US,
  tagging=on,
  tagging-setup={math/setup=mathml-SE,tabsorder=structure}
}
\documentclass[11pt,letterpaper]{article}
\usepackage[margin=0.68in,includefoot]{geometry}
\usepackage{newcomputermodern}
\usepackage{amsmath,amscd,mathtools}
\usepackage{tikz,adjustbox}
\providecommand{\square}{\mdlgwhtsquare}
\usepackage[hidelinks]{hyperref}
\hypersetup{
  pdftitle={Complex Analysis PhD Exam, August 2025},
  pdfauthor={University of Florida Department of Mathematics},
  pdfsubject={Graduate mathematics examination},
  pdfdisplaydoctitle=true,
  bookmarksopen=true
}
\setcounter{secnumdepth}{0}
\setlength{\parindent}{0pt}
\setlength{\parskip}{0.28em}
\setlength{\emergencystretch}{3em}
\setlength{\leftmargini}{2.15em}
\setlength{\leftmarginii}{2.65em}
\setlength{\leftmarginiii}{3em}
\renewcommand{\labelenumi}{\textbf{\theenumi.}}
\pagestyle{empty}
\raggedbottom
\allowdisplaybreaks
\newenvironment{examproblems}[1]{%
  \begin{enumerate}%
  \setcounter{enumi}{#1}%
  \setlength{\topsep}{0.35em}%
  \setlength{\partopsep}{0pt}%
  \setlength{\itemsep}{0.48em}%
  \setlength{\parsep}{0.18em}%
}{\end{enumerate}}
\newenvironment{examsubparts}{%
  \begin{enumerate}%
  \setlength{\topsep}{0.18em}%
  \setlength{\partopsep}{0pt}%
  \setlength{\itemsep}{0.16em}%
  \setlength{\parsep}{0.1em}%
}{\end{enumerate}}
\newenvironment{examinstructions}{%
  \par\begingroup\itshape
}{%
  \par\endgroup\vspace{0.18em}
}
\makeatletter
\renewcommand\section{\@startsection{section}{1}{0pt}%
  {-1.25ex plus -.25ex minus -.1ex}%
  {0.55ex plus .1ex}%
  {\normalfont\large\bfseries}}
\renewcommand{\@maketitle}{%
  \newpage\null\vskip -1em%
  {\footnotesize\bfseries
    \href{https://www.ufl.edu/}{University of Florida}, \href{https://math.ufl.edu/}{Department of Mathematics}\par}%
  \vskip -0.5em%
  \tagstructbegin{tag=Title}%
    {\LARGE\bfseries\href{https://gma.math.ufl.edu/exams/complex-analysis-phd/}{Complex Analysis PhD Exam}, August 2025\par}%
  \tagstructend%
  \vskip 0.25em%
  \tagmcbegin{artifact}%
    \hrule height 1pt%
  \tagmcend%
  \vskip 1em%
}
\makeatother
\title{Complex Analysis PhD Exam, August 2025}
\author{}
\date{}
\begin{document}
\maketitle
\thispagestyle{empty}
\begin{examinstructions}
Answer SIX of the following eight questions. Explain your work in a precise and logical way. Any standard result may be used without proof but must be stated clearly and correctly. Throughout the exam, \(\mathbb{D}\) is the open unit disk in~\(\mathbb{C}\).
\end{examinstructions}
\begin{examproblems}{0}
\item
Let~\(f\) be an analytic function on some open connected domain~\(G\) and suppose that there exists an \(a\in\mathbb{C}\) such that \(f(z)=a\) has uncountably many solutions in~\(G\). What can be said about~\(f\)?
\item
State Morera’s theorem and use it to prove that if \(\{f_n\}\subset H(\mathbb{D})\) is a sequence of analytic functions that converges locally uniformly on~\(\mathbb{D}\) to a function~\(f\), then~\(f\) is analytic.
\item
Determine all entire functions~\(f\) that satisfy \(f(f(z))=z\).
\item
Let \(p(z)=z^4+3z^3-1\). Is \(1/p(z)\) analytic on the annulus \(1<|z|<4\)?
\item
Evaluate the following improper integral: 
\[
\int_0^\infty \frac{1}{1+x^4}\,dx
\]
\item
Let~\(f\) be a nonconstant entire function that is periodic with nonzero period \(0\ne T\in\mathbb{C}\). Prove that \(f(z)-z=0\) has infinitely many solutions.
\item
Determine an explicit analytic bijection from \(\mathbb{D}\cap\{z:\operatorname{im}(z)>0\}\) to~\(\mathbb{D}\).
\item
Let~\(G\) be an open domain and let \(\{f_n\}\subset H(G)\). Also let~\(f\) be a nonconstant function such that \(f_n\to f\) in \(H(G)\). Let \(a\in G\) and set \(\alpha=f(a)\). Show that there is a sequence \(\{a_n\}\subset G\) such that \(a_n\to a\) and \(f_n(a_n)=\alpha\) for all~\(n\) sufficiently large.
\end{examproblems}
\end{document}
